\documentclass[10pt,a4paper]{amsart}
\usepackage[margin=19mm]{geometry}
\usepackage{amsmath,amssymb,mathtools}
\usepackage[scr=rsfs,cal=euler]{mathalfa}
\usepackage{xfrac}
\usepackage[unicode,hidelinks,bookmarks=false]{hyperref}
\newcommand{\ie}{i.e.\ }
\newcommand{\cf}{cf.\ }
\newcommand{\resp}{respectively\ }
\newcommand{\Subsection}{\subsection}
\providecommand{\nobreakdash}{\nobreak\hbox{-}}
\setlength{\emergencystretch}{2em}
\multlinegap0pt
\usepackage{amssymb}
\usepackage{mathtools}
\newtheorem{lemma}{Lemma}
\newtheorem{theorem}{Theorem}
\newcommand{\C}{\mathbb{C}}
\newcommand{\R}{\mathbb{R}}
\title{The $L^2$ Norm of the Interior Cauchy Transform: Beyond the First Dirichlet Eigenvalue}
\author{David Kalaj}
\date{Source: arXiv:2602.13740v1 (CC BY 4.0)}
\begin{document}
\maketitle

\noindent\emph{Source and licence.} The $L^2$ Norm of the Interior Cauchy Transform: Beyond the First Dirichlet Eigenvalue. Version arXiv:2602.13740v1 (CC BY 4.0), distributed by arXiv under the Creative Commons Attribution 4.0 International licence, http://creativecommons.org/licenses/by/4.0/, as recorded in the arXiv licence metadata for this exact version and shown on the abstract page https://arxiv.org/abs/2602.13740v1. Full licence notice: \texttt{licenses/arxiv-2602.13740-CC-BY-4.0.txt}.\par\smallskip

\noindent\textit{Editorial scope.} Counted unit: the article's theorem thm:rigidity-disk (source0.tex lines 603--618). The article's complete original proof of it is delivered with it (source0.tex lines 620--713, one \texttt{proof} environment of 94 lines). No mathematics is written, repaired or re-derived: the statement and the proof are the article's own lines, in the article's own notation, and no retained line is substituted or re-flowed.  Retained uncounted as the source-local context the proof needs: lines 511--523.  Nothing was omitted from any retained span, so the package carries no omission record.  No retained line contains a citation, so the wrapper carries no bibliography and no bibliography entry.  No retained line contains a figure environment, an image-inclusion command or drawing code, so no image is delivered and the package declares no asset dependency.  Editorial formatting only, in addition: an AMS \texttt{amsart} preamble that restates the article's own declarations and macros used by the retained lines, namely the environments \texttt{lemma}, \texttt{theorem} on separate counters, together with the standard packages those lines need.  The retained environments print in this document as Lemma 1, Theorem 1 in order of appearance, whereas the article prints the same items with the numbering of its own full document.  The complete native source of the article as distributed in the arXiv e-print for this exact version is bundled unchanged as \texttt{sources/194663/source0.tex}.
\begin{lemma}[Exterior holomorphy and decay of the Cauchy transform]\label{lem:ext-holo-decay}
Let $D\subset\C$ be bounded and let $f\in L^1(D)$. Define
\[
F(z):=\frac1\pi\int_D \frac{f(w)}{z-w}\,dA(w),\qquad z\in\Omega:=\C\setminus\overline D.
\]
Then $F$ is holomorphic on $\Omega$ and, as $z\to\infty$,
\[
F(z)=\frac{1}{\pi}\left(\frac{M_0}{z}+\frac{M_1}{z^2}+\frac{M_2}{z^3}+\cdots\right),
\qquad
M_k:=\int_D f(w)\,w^k\,dA(w),
\]
in particular $F(z)=O(1/z)$.
\end{lemma}

\begin{theorem}[Rigidity]\label{thm:rigidity-disk}
Let $D \subset \C$ be a bounded, simply connected domain with $C^{1,\alpha}$ boundary.
Let $u$ be the first Dirichlet eigenfunction,
\[
-\Delta u=\lambda_1(D)\,u \quad \text{in } D, \qquad u=0 \quad \text{on } \partial D,
\]
and define
\[
v_0 := -\frac{4}{\lambda_1(D)}\,\partial_z u .
\]
Then
\[
C_D u \equiv v_0 \quad \text{in } D
\]
if and only if $D$ is a disk.
\end{theorem}

\begin{proof}
Assume first that $C_Du\equiv v_0$ in $D$, where $v_0=-(4/\lambda_1)\partial_z u$.
Let $\Omega:=\C\setminus\overline D$ and consider the exterior Cauchy transform
\[
F(z):=\frac{1}{\pi}\int_D \frac{u(w)}{z-w}\,dA(w),\qquad z\in\Omega.
\]
By Lemma~\ref{lem:ext-holo-decay}, $F$ is holomorphic on $\Omega$ and satisfies $F(z)=O(1/z)$ as $z\to\infty$.

For smooth $u$ and $\partial D$, the same integral defines a continuous function on $\C$, and its restriction
to $D$ coincides with $C_Du$. Hence the identity $C_Du=v_0$ on $D$ implies that $v_0$ has a holomorphic
extension to $\Omega$, given by
\[
\partial_z u(w) := -\frac{\lambda_1}{4}\,F(w),\qquad w\in\Omega.
\]
In particular, this extension is holomorphic on $\Omega$ and
\begin{equation}\label{eq:dz-decay}
\partial_z u(w)=O(1/w)\qquad (w\to\infty).
\end{equation}

\paragraph{Boundary identity.}
On $\partial D$ we have $u=0$, so $\nabla u$ is normal:
\[
\nabla u = (\partial_n u)\,n,
\]
where $n=n_x+i n_y$ is the unit outward normal. Using $\partial_z=\tfrac12(\partial_x-i\partial_y)$, we obtain
\[
\partial_z u
=\frac12(\partial_n u)\,\overline{n}
\qquad \text{on }\partial D.
\]
Let $\gamma=\gamma(s)$ be an arclength parametrization of $\partial D$. Then $|\gamma'(s)|=1$ and
\[
n(s)=-i\,\gamma'(s),
\qquad\text{hence}\qquad
\overline{n(s)}=i\,\overline{\gamma'(s)}=\frac{i}{\gamma'(s)}.
\]
Consequently,
\[
\partial_z u(\gamma(s))\,d\gamma
=\partial_z u(\gamma(s))\,\gamma'(s)\,ds
=\frac12(\partial_n u)\,\overline{n}\,\gamma'(s)\,ds
=\frac{i}{2}(\partial_n u)\,ds,
\]
so the complex differential $\partial_z u\,d\gamma$ has vanishing real part on $\partial D$
(equivalently, it is purely imaginary along $\partial D$).

\paragraph{Rigidity via the exterior conformal map.}
Let $\Phi:\Omega\to\{|\zeta|>1\}$ be the exterior conformal map normalized by $\Phi(\infty)=\infty$,
and write
\[
Z(\zeta):=\Phi^{-1}(\zeta).
\]
Define
\[
\omega(\zeta):=\bigl(\partial_z u\bigr)\bigl(Z(\zeta)\bigr)\,Z'(\zeta),
\qquad |\zeta|>1.
\]
Then $\omega$ is holomorphic for $|\zeta|>1$, and \eqref{eq:dz-decay} together with the normalization
\[
Z(\zeta)=a\zeta+b+O(1/\zeta)\qquad (\zeta\to\infty)
\]
implies
\[
\omega(\zeta)=O(1/\zeta)\qquad (\zeta\to\infty),
\]
so $F_1(\zeta):=\zeta\,\omega(\zeta)$ is holomorphic on $|\zeta|>1$ and bounded at $\infty$.

On $|\zeta|=1$ we have $Z(\zeta)\in\partial D$ and $dZ=Z'(\zeta)\,d\zeta$, hence
\[
\omega(\zeta)\,d\zeta
=\bigl(\partial_z u\bigr)\bigl(Z(\zeta)\bigr)\,dZ.
\]
The boundary identity above shows that this differential has vanishing real part on $|\zeta|=1$; therefore,
for $\zeta=e^{it}$ (so $d\zeta=i\zeta\,dt$) we obtain
\[
\Re\!\bigl(F_1(e^{it})\,i\,dt\bigr)=0
\qquad\Longrightarrow\qquad
F_1(e^{it})\in\R,\quad t\in\R.
\]
Thus $F_1$ extends holomorphically across $|\zeta|=1$ by Schwarz reflection and is an entire bounded
function; hence $F_1$ is constant, say $F_1(\zeta)\equiv c\in\R$, and thus
\[
\omega(\zeta)=\frac{c}{\zeta}.
\]
Unwinding the definition of $\omega$ shows that $|Z'(\zeta)|$ is constant on $|\zeta|=1$.
Since $\log|Z'(\zeta)|$ is harmonic on $|\zeta|>1$, constant boundary values imply $\log|Z'|$ is constant
on $|\zeta|>1$, so $Z'(\zeta)$ is constant and $Z(\zeta)$ is a M\"obius map. Consequently $\partial D$ is a circle
and $D$ is a disk.


\noindent\textbf{Conversely.}
If $D$ is a disk, the first eigenfunction is radial and a direct computation (e.g.\ in polar coordinates)
shows that $C_D u = -\frac{4}{\lambda_1(D)}\,\partial_z u$ holds identically.
\end{proof}

\end{document}
